\chapter{Complexity — Clean Sweep}
%%% generated by the blueprint pipeline from TCSlib.Complexity.SpaceComplexity.Machines.CleanSweep
\section{Overview}

Declarations of the Lean module \texttt{TCSlib.Complexity.SpaceComplexity.Machines.CleanSweep}.

\section{Declarations}

\begin{definition}[Cleaning phases]
\lean{Complexity.LogProg.CPh}
\label{Complexity.LogProg.CPh}
\leanok
The four cleaning phases are marking the current cell, moving right over marked cells,
moving left while erasing marked cells, and moving right back to the origin mark while
erasing it.
\end{definition}

\begin{definition}[Cleaning machine states]
\lean{Complexity.LogProg.CleanSt}
\label{Complexity.LogProg.CleanSt}
\leanok
\uses{Complexity.LogProg.CPh}
For a state type $S$ and a natural number $k$ of work tapes, the cleaning machine has
states for initialization from each state $q\in S$, simulation from each state $q\in S$,
and cleaning each tape $i<k$ in any of the four cleaning phases.
\end{definition}

\begin{definition}[Idle work-tape action]
\lean{Complexity.LogProg.idleK}
\label{Complexity.LogProg.idleK}
\leanok
For a natural number $k$, the idle action on each of the $k$ work tapes is the pair
consisting of no symbol and the sign $0$.
\end{definition}

\begin{definition}[Cleaning-phase tape action]
\lean{Complexity.LogProg.clAct}
\label{Complexity.LogProg.clAct}
\leanok
\uses{Complexity.LogProg.CleanSt, Turing.Action}
For a cleaning tape index $i$ among $k$ tapes, an optional work-tape write instruction
$w_D$, an optional mark-tape write instruction $w_M$, a head movement $m\in\{-1,0,1\}$,
and an optional successor cleaning state $s$, define an action on $2k$ work tapes that
leaves the input head and output unchanged. On tape $i$ in each block of $k$ tapes it
performs the respective write instruction and movement $m$; on every other tape it
leaves the symbol and head position unchanged. Its successor state is $s$.
\end{definition}

\begin{definition}[Next cleaning tape]
\lean{Complexity.LogProg.clNext}
\label{Complexity.LogProg.clNext}
\leanok
\uses{Complexity.LogProg.CleanSt}
For an index $i$ among $k$ tapes, define the next cleaning state to be the cleaning
state for tape $i+1$ in the marking phase when $i+1<k$, and undefined otherwise.
\end{definition}

\begin{definition}[Cleaned machine]
\lean{Complexity.LogProg.cleanTM}
\label{Complexity.LogProg.cleanTM}
\leanok
\uses{Complexity.LogProg.CleanSt, Complexity.LogProg.clAct, Complexity.LogProg.clNext, Complexity.LogProg.idleK, Turing.MultiTapeTM}
Given a deterministic Turing machine $M$ with $k$ work tapes, Boolean alphabet, and
state set $S$, and an initial state $q_0\in S$, define a machine with $2k$ work tapes
and cleaning states. It starts in the initialization state for $q_0$, marks the origin
on the added tapes, and then simulates $M$ on the first $k$ tapes while maintaining
marks on the added tapes; when $M$ halts, it proceeds through the cleaning phases to
erase those marks and return the work tapes to blank, halting after the cleaning
process.
\end{definition}

\begin{definition}[Combined tape configuration]
\lean{Complexity.LogProg.ccfg}
\label{Complexity.LogProg.ccfg}
\leanok
\uses{Complexity.LogProg.CleanSt, Turing.Cfg}
For a Boolean input string $x$, an optional cleaning state $s$, an input-head position,
two collections of $k$ work-tape contents $D_t,M_t$ and head positions $D_p,M_p$, and an
output string, define the configuration on $2k$ work tapes by concatenating the two
collections of tape contents and the two collections of head positions, in that order.
\end{definition}

\begin{lemma}[State of the Combined Configuration]
\lean{Complexity.LogProg.ccfg_state}
\label{Complexity.LogProg.ccfg_state}
\leanok
\uses{Complexity.LogProg.CleanSt, Complexity.LogProg.ccfg}
\difficulty{1}
Let $x$ be a finite Boolean input, let $k$ be a natural number, and choose an optional
control state, an input-head position, two families of $k$ work tapes with their head
positions, and an output word. The configuration formed by concatenating the two tape
families has the chosen optional control state.
\end{lemma}

\begin{lemma}[Marked Tape Reading]
\lean{Complexity.LogProg.ccfg_mark}
\label{Complexity.LogProg.ccfg_mark}
\leanok
\uses{Complexity.LogProg.CleanSt, Complexity.LogProg.ccfg, Turing.Cfg.workTapeSymbols}
\difficulty{1}
Let $x$ be a finite Boolean input, let $k$ be a natural number, and let the control
state, input-head position, two families of $k$ work tapes with their head positions,
and output word be arbitrary. Form a configuration with $2k$ work tapes by placing the
first family before the second. For each $i<k$, the symbol read from tape $k+i$ is the
content of the $i$th tape in the second family at its head position.
\end{lemma}

\begin{lemma}[Reading the First Tape Family]
\lean{Complexity.LogProg.ccfg_work}
\label{Complexity.LogProg.ccfg_work}
\leanok
\uses{Complexity.LogProg.CleanSt, Complexity.LogProg.ccfg, Turing.Cfg.workTapeSymbols}
\difficulty{1}
Let $x$ be a finite Boolean input, let $k$ be a natural number, and choose an optional
control state, an input-head position, two families of $k$ work tapes with their head
positions, and an output word. For each $i<k$, the symbol read from tape $i$ in the
combined configuration is the content of the $i$th tape in the first family at its head
position.
\end{lemma}

\begin{definition}[Tape write update]
\lean{Complexity.LogProg.wr}
\label{Complexity.LogProg.wr}
\leanok
Given an optional write instruction $w$ and a tape-content function $f:\mathbb
Z\to\mathrm{Option}(\mathrm{Bool})$, define the updated tape at position $p$ to be $f$
unchanged when $w$ is absent; when $w$ is present with value $s$, replace the content at
$p$ by $s$ and leave all other positions unchanged.
\end{definition}

\begin{lemma}[Effect of a paired-tape cleaning action]
\lean{Complexity.LogProg.apply_clAct}
\label{Complexity.LogProg.apply_clAct}
\leanok
\uses{Complexity.LogProg.CleanSt, Complexity.LogProg.ccfg, Complexity.LogProg.clAct, Complexity.LogProg.wr, Turing.Action.apply, Turing.moveInputPos_zero}
\difficulty{5}
Let the two tape banks each have $k$ tapes, and let a configuration have input $x$,
state $st$, input-head position $ip$, tape contents $D_t,M_t$, head positions $D_p,M_p$,
and output $out$. For a tape index $i$, write instructions $w_D,w_M$, movement sign
$mv$, and successor state $st'$, applying the cleaning action gives state $st'$, the
same input-head position and output, and updates only tape $i$ in each bank: its
contents become the result of applying the respective write instruction at its current
head position, and its head position increases by $mv$.
\end{lemma}

\begin{definition}[Interval marking function]
\lean{Complexity.LogProg.markI}
\label{Complexity.LogProg.markI}
\leanok
For integers $a,b$, define a tape-content function on integer positions $z$ that stores
true at the origin $z=0$, stores false at positions satisfying $a\le z\le b$, and is
blank at all other positions.
\end{definition}

\begin{definition}[Work-tape position bound]
\lean{Complexity.LogProg.PB}
\label{Complexity.LogProg.PB}
\leanok
\uses{Complexity.LogProg.CleanSt, Turing.Cfg}
For a configuration $g$ of the cleaning machine with $2k$ work tapes and Boolean input
$x$, define the property that every work-tape head position has absolute value at most
$B$.
\end{definition}

\begin{lemma}[Head bounds for the combined configuration]
\lean{Complexity.LogProg.PB_ccfg}
\label{Complexity.LogProg.PB_ccfg}
\leanok
\uses{Complexity.LogProg.CleanSt, Complexity.LogProg.PB, Complexity.LogProg.ccfg}
\difficulty{4}
Let $B\in\mathbb{Z}$, and form a configuration with $2k$ work tapes by concatenating the
contents and head positions of two $k$-tape collections, with any state, input-head
position, and output. If every head position in each collection has absolute value at
most $B$, then every work-tape head position in the combined configuration has absolute
value at most $B$.
\end{lemma}

\begin{lemma}[Absolute-Value Bound Under Update]
\lean{Complexity.LogProg.abs_update_le}
\label{Complexity.LogProg.abs_update_le}
\leanok
\difficulty{4}
Let $k$ be a natural number, let $D$ assign an integer to each of the $k$ indices, and
let $i$ be one such index. If every value of $D$ has absolute value at most $B$, and an
integer $v$ also has absolute value at most $B$, then the function obtained by replacing
$D(i)$ by $v$ has absolute value at most $B$ at every index.
\end{lemma}

\begin{lemma}[Rightward cleanup run]
\lean{Complexity.LogProg.goR_run}
\label{Complexity.LogProg.goR_run}
\leanok
\uses{Complexity.LogProg.PB, Complexity.LogProg.PB_ccfg, Complexity.LogProg.abs_update_le, Complexity.LogProg.apply_clAct, Complexity.LogProg.ccfg, Complexity.LogProg.ccfg_mark, Complexity.LogProg.ccfg_state, Complexity.LogProg.cleanTM, Complexity.LogProg.markI, Complexity.LogProg.wr, Turing.MultiTapeTM.runFrom, Turing.MultiTapeTM.runFrom_succ_eq_step, Turing.MultiTapeTM.step}
\difficulty{5}
Let $M$ be a $k$-work-tape Turing machine, let $q_0$ be an initial state, and fix an
input, an index $i<k$, an input-head position, output word, and integer bound $B$. Let
$a,b$ be integers, and let $D_t,M_t$ be tape contents and $D_p,M_p$ be head positions
for the two banks of work tapes. Suppose $M_t(i)$ is the tape marking the origin and the
interval from $a$ through $b$, $0\le b$, and $b+1\le B$. For every $n$ and head-position
function $D_p$ such that $b+1-D_p(i)=n$, $a\le D_p(i)$, and $|D_p(j)|\le B$ for every
tape index $j$, the cleanup machine started in its rightward-walk phase at index $i$ has
all work-head positions bounded in absolute value by $B$ at every time $t\le n+1$. At
time $n+1$ it is in the erasing phase at index $i$, with both banks' head positions
updated at $i$ to $b$ and all other configuration data unchanged.
\end{lemma}

\begin{definition}[Erase above a position]
\lean{Complexity.LogProg.eraseAbove}
\label{Complexity.LogProg.eraseAbove}
\leanok
Given a tape-content function $f:\mathbb Z\to\mathrm{Option}(\mathrm{Bool})$ and an
integer position $p$, define a new tape that retains $f(z)$ for every $z\le p$ and is
blank at every $z>p$.
\end{definition}

\begin{lemma}[Tape Erasure Run and Head Bound]
\lean{Complexity.LogProg.erase_run}
\label{Complexity.LogProg.erase_run}
\leanok
\uses{Complexity.LogProg.PB, Complexity.LogProg.PB_ccfg, Complexity.LogProg.abs_update_le, Complexity.LogProg.apply_clAct, Complexity.LogProg.ccfg, Complexity.LogProg.ccfg_mark, Complexity.LogProg.ccfg_state, Complexity.LogProg.cleanTM, Complexity.LogProg.eraseAbove, Complexity.LogProg.markI, Complexity.LogProg.wr, Turing.MultiTapeTM.runFrom, Turing.MultiTapeTM.runFrom_succ_eq_step, Turing.MultiTapeTM.step}
\difficulty{6}
Let $M$ be a Boolean multi-tape Turing machine with initial state $q_0$, let $i$ be a
work-tape index, let the input-head position be $ip$, and let the output word be $out$.
For integers $a,b,B$ with $a\le 0$ and $-B\le a-1$, a tape content function $f$, and any
$n\in\mathbb N$, suppose the data tapes $D_t$, their head positions $D_p$, and the mark
tapes $M_t$ satisfy $D_p(i)-(a-1)=n$, $D_p(i)\le b$, $D_t(i)$ agrees with $f$ on cells
at most $D_p(i)$ and is blank above them, $M_t(i)$ has the prescribed origin and
interval marks from $a$ through $D_p(i)$, and every data-tape head position has absolute
value at most $B$. Starting from the cleaning machine's erase phase at index $i$ with
these tapes and head positions, every configuration through time $n+1$ has all work-tape
head positions bounded in absolute value by $B$. At time $n+1$, the machine is in the
back phase at index $i$; the selected data tape contains $f$ through cell $a-1$ and is
blank above, its head is at $a$, and its mark tape has the prescribed marks through cell
$a-1$ with head at $a$; all other tapes, the input-head position, and the output are
unchanged.
\end{lemma}

\begin{lemma}[Marking an empty integer interval]
\lean{Complexity.LogProg.markI_empty}
\label{Complexity.LogProg.markI_empty}
\leanok
\uses{Complexity.LogProg.markI}
\difficulty{3}
For all integers $a$ and $z$, marking $z$ relative to the interval from $a$ to $a-1$
returns true exactly when $z=0$, and is undefined otherwise.
\end{lemma}

\begin{lemma}[Bounds during a tape-cleaning return phase]
\lean{Complexity.LogProg.back_run}
\label{Complexity.LogProg.back_run}
\leanok
\uses{Complexity.LogProg.PB, Complexity.LogProg.PB_ccfg, Complexity.LogProg.abs_update_le, Complexity.LogProg.apply_clAct, Complexity.LogProg.ccfg, Complexity.LogProg.ccfg_mark, Complexity.LogProg.ccfg_state, Complexity.LogProg.clNext, Complexity.LogProg.cleanTM, Complexity.LogProg.markI, Complexity.LogProg.markI_empty, Complexity.LogProg.wr, Turing.MultiTapeTM.runFrom, Turing.MultiTapeTM.runFrom_succ_eq_step, Turing.MultiTapeTM.step}
\difficulty{5}
Let $M$ be a $k$-tape Boolean Turing machine, let $q_0$ be a state, and fix an input, a
tape index $i$, an input-head position, an initial output, and an integer bound $B$. For
an integer $a$, suppose there are a natural number $n$ and tape contents and head
positions $D_t,M_t,D_p$ such that $-D_p(i)=n$, $a\le D_p(i)$, tape $M_t(i)$ has its
origin mark at $0$ and is blank outside $[a,a-1]$, and every initial head position in
$D_p$ has absolute value at most $B$. Starting the cleaned machine in cleaning phase
\emph{back} for tape $i$, every configuration through time $n+1$ has every work-tape
head position of absolute value at most $B$. At time $n+1$, its state is the next
cleaning phase for tape $i$ (or halted if there is no next tape), its first bank of tape
contents is unchanged, and both banks' head positions at tape $i$ are reset to $0$ while
the second bank's tape $i$ is blank; all other tape contents and positions and the input
head and output are unchanged.
\end{lemma}

\begin{lemma}[Cleaning One Tape]
\lean{Complexity.LogProg.clean_tape}
\label{Complexity.LogProg.clean_tape}
\leanok
\uses{Complexity.LogProg.PB, Complexity.LogProg.PB_ccfg, Complexity.LogProg.abs_update_le, Complexity.LogProg.apply_clAct, Complexity.LogProg.back_run, Complexity.LogProg.ccfg, Complexity.LogProg.ccfg_mark, Complexity.LogProg.ccfg_state, Complexity.LogProg.clNext, Complexity.LogProg.cleanTM, Complexity.LogProg.eraseAbove, Complexity.LogProg.erase_run, Complexity.LogProg.goR_run, Complexity.LogProg.markI, Complexity.LogProg.runFrom_forall_append, Complexity.LogProg.wr, Turing.MultiTapeTM.runFrom, Turing.MultiTapeTM.runFrom_add, Turing.MultiTapeTM.step}
\difficulty{6}
Let $M$ be a $k$-tape Turing machine with Boolean symbols, choose its initial state
$q_0$, a Boolean input $x$, a tape index $i<k$, an input-head position, an output word,
an integer bound $B$, integer endpoints $a\le b$, and tape contents $D_t,M_t$ with head
positions $D_p$. Suppose $a\le0\le b$, $-B\le a-1$, $b+1\le B$, and $a\le D_p(i)\le b$.
Suppose that writing false at the scanned cell of the $i$th tape of $M_t$ when it does
not already contain true, and otherwise leaving it unchanged, produces the tape that has
true at the origin, false on the interval from $a$ to $b$, and blank elsewhere. Also
suppose the nonblank cells of $D_t(i)$ lie in $[a,b]$, and every head position in $D_p$
has absolute value at most $B$. Then the cleaning machine, started in the cleaning phase
for tape $i$ at its marking stage with these tapes and positions, reaches some time $T$
by which every work-tape head position at every time $t\le T$ has absolute value at most
$B$, and its configuration at time $T$ has the next cleaning state, unchanged input-head
position and output, and both tape families blank with all their heads at zero except
that the stated tape index advances according to the cleaning successor.
\end{lemma}
